\scanpage{3}
\mspage{2}
contexte d'un atelier définissant une ``contrée''.

Le travail précédent a permis de dégager le point de vue le plus commode pour le formalisme algébrique. La notion fondamentale me semble bien celle de \uline{figure}, et non pas celle de ``figure élémentaire'' ou ``multistrate'', voire celle de ``lieu'' (voire, la notion de ``lieu ponctuel'' ou ``point'', dont il n'a pas encore été question\ldots). Une des raisons en est la simplicité du formalisme algébrique,
% slanted
\margin{L'autre raison, c'est que, dans le cas d'une ``application $f$'' d'une ``région'' dans une autre, l'image inverse d'une ``figure'' (\unsure{quand|quant} \unsure{on} peut donner un sens \unsure{géométrique} à ce \unsure{mot|terme}) \deleted{\unsure{sera}} tend à ne \unsure{pas} être \unsure{encore} une figure, mais l'image réciproque d'une multistr. est une multistr.}%
 où il est agréable de pouvoir manipuler les lieux et multistrates comme étant des figures particulières. Quand une figure \inserted{$F$} est interprétée \inserted{(par l'ens de ses ``strates'')} comme partie de l'ens $\unsure{M|\mathcal M}$ \inserted{(ordonné par $\leq$)} des multistrates, nous allons écrire $\widetilde F$ (et non $F$) pour cette partie. Quand elle est interprétée comme une figure \inserted{ensembliste} dans $M$, nous l'écrivons $\mathrm{Mulomb}(F)$ (``grande multiombre''), son support est $\mathrm{Omb}(F)$ (la ``grande ombre'' de $F$) -- cette dernière, contrairement à $\widetilde F$ et à $\mathrm{Mulomb}(F)$, ne permet pratiquement jamais de réconstituer $F$, \deleted{\unsure{le|il}} mais moralement, son ``support'' seulement
